% The pilot, the arithmetic that explained its null, and the long matrix
% that retracted the explanation while keeping the null. The body states
% what the two matrices jointly establish; this is how they got there.
\paragraph{Result: no learning-curve effect, on either horizon.}
The pilot was five arms over the two shortest sequences at three seeds:
615 tasks, every one scored by the official docker harness, no failed
cells. The pre-registered comparison did not hold, and the arm carrying
no store at all had the highest resolve rate. We had an explanation---at
$0.05$ upkeep against $1.0$ spawn energy an entry starves after exactly
$20$ cycles, these sequences are $19$ and $22$ tasks, and with nothing
starving the ledger \emph{is} \texttt{keep\_everything} plus
settlement---and rather than rest on it we tested it.
\pointer{\S\ref{sec:swebench-pilot} gives the pilot matrix and that
arithmetic in full.}
\inplace{experiments-swebench-pilot}

\paragraph{The long matrix: the excuse does not survive its own test.}
We re-ran the full design over \texttt{django} and \texttt{sympy}, $50$
tasks each at $2.5\times$ the starvation horizon: another 30 cells,
$1{,}500$ tasks, no failed cells (Table~\ref{tab:swebench-long}). The
mechanism half of the prediction held exactly. Starvation moved from
$4\%$ of minted lessons to $52\%$---$155$ of $300$ under
\texttt{memory\_on}, at a steady $21$--$23$ per cell---and the surviving
population held at $24$ where \texttt{keep\_everything} climbs to the
full $50$. Leanness, the one property that survived every earlier
comparison, is demonstrated on real tasks here for the first time: half
the store, and the deaths begin at position $20$ where the arithmetic
says they must.

The capability half did not hold. With removal fully active the
pre-registered comparison is $+0.027$ $[-0.020, +0.073]$, permutation
$p = 0.50$; on \texttt{django} alone, where per-seed variance is
smallest ($23,23,23$ for \texttt{memory\_off}), it is $-0.013$ with
$p = 1.0000$. Halving the store changed measured capability by nothing,
and \texttt{memory\_off} is again the highest-scoring arm ($0.360$).
Across all four sequences and $2{,}115$ evaluated tasks, no memory arm
has beaten no memory at all.

We therefore retract the benchmark explanation as a defence of the
mechanism and keep it only as a description of the pilot. What the two
matrices jointly establish is narrower than what we set out to show and
more useful than a second ambiguous null: on real software-engineering
tasks outcome-based selection produces a smaller store; the absence of a
detected capability difference does not establish preserved success, and buys no measured capability either.


